\documentclass{article}
\usepackage{amsmath,amssymb,amsthm}
\usepackage{algorithm}
\usepackage{algpseudocode}
\usepackage{enumitem}
\usepackage{hyperref}
\usepackage{geometry}
\geometry{margin=1in}
\newtheorem{theorem}{Theorem}
\newtheorem{lemma}{Lemma}
\newtheorem{assumption}{Assumption}
\newcommand{\E}{\mathbb{E}}
\newcommand{\prox}{\operatorname{prox}}
\newcommand{\norm}[1]{\left\lVert#1\right\rVert}
\begin{document}

\section*{Algorithm Name}
\textbf{Stochastic Variance‑Reduced Gradient (SVRG) with Double‑Loop Structure}

\section*{Mathematical Formulation}
Let the finite‑sum objective be 
\[
f(x)=\frac{1}{n}\sum_{i=1}^{n} f_i(x),\qquad f_i:\mathbb{R}^d\to\mathbb{R}.
\]
Denote by $x^\star$ the unique minimizer of $f$.  
SVRG proceeds in \emph{epochs} $s=0,1,\dots,S-1$.  
At the beginning of epoch $s$ a \emph{snapshot} $\tilde x^{\,s}$ is fixed and the \emph{full gradient}
\[
\tilde \mu^{\,s}= \nabla f(\tilde x^{\,s})=\frac{1}{n}\sum_{i=1}^{n}\nabla f_i(\tilde x^{\,s})
\tag{1}
\]
is computed once.

Inside epoch $s$ we perform $m$ inner updates indexed by $t=0,\dots,m-1$:
\[
\begin{aligned}
\text{Sample } i_t &\sim \text{Uniform}\{1,\dots,n\},\\[0.2em]
v_t^{\,s} &= \nabla f_{i_t}\bigl(x_t^{\,s}\bigr)-\nabla f_{i_t}\bigl(\tilde x^{\,s}\bigr)+\tilde\mu^{\,s},\\[0.2em]
x_{t+1}^{\,s} &= x_t^{\,s}-\eta\, v_t^{\,s},
\end{aligned}
\tag{2}
\]
where $\eta>0$ is a constant stepsize.  
After the $m$ inner steps we set the next snapshot as the \emph{averaged} inner iterate
\[
\tilde x^{\,s+1}= \frac{1}{m}\sum_{t=1}^{m} x_t^{\,s}.
\tag{3}
\]

The algorithm terminates after $S$ epochs and returns $\tilde x^{\,S}$.

\section*{Pseudocode}
\begin{algorithm}[H]
\caption{SVRG (double‑loop)}\label{alg:svrg}
\begin{algorithmic}[1]
\Require Initial point $\tilde x^{\,0}$, stepsize $\eta$, inner length $m$, epochs $S$
\For{$s=0$ \textbf{to} $S-1$}
    \State Compute full gradient $\tilde\mu^{\,s}= \nabla f(\tilde x^{\,s})$ \Comment{Eq.~(1)}
    \State $x_0^{\,s}\gets \tilde x^{\,s}$
    \For{$t=0$ \textbf{to} $m-1$}
        \State Sample $i_t\stackrel{\text{i.i.d.}}{\sim}\text{Unif}\{1,\dots,n\}$
        \State $v_t^{\,s}= \nabla f_{i_t}(x_t^{\,s})-\nabla f_{i_t}(\tilde x^{\,s})+\tilde\mu^{\,s}$ \Comment{Eq.~(2)}
        \State $x_{t+1}^{\,s}= x_t^{\,s}-\eta\,v_t^{\,s}$ \Comment{Eq.~(2)}
    \EndFor
    \State $\tilde x^{\,s+1}= \frac{1}{m}\sum_{t=1}^{m} x_t^{\,s}$ \Comment{Eq.~(3)}
\EndFor
\State \Return $\tilde x^{\,S}$
\end{algorithmic}
\end{algorithm}

\section*{Dry Run Example}
Consider a single‑sample problem ($n=1$) with
\[
f(w) = (w-3)^2,
\qquad \nabla f(w)=2(w-3).
\]
Take $\tilde x^{\,0}=0$, stepsize $\eta=0.1$, inner length $m=3$, and run one epoch ($S=1$).

\begin{center}
\begin{tabular}{c|c|c|c}
$t$ & $x_t^{\,0}$ & $v_t^{\,0}$ (variance‑reduced) & $x_{t+1}^{\,0}=x_t^{\,0}-\eta v_t^{\,0}$\\ \hline
0 & $0$ &
\begin{aligned}
&\nabla f_{1}(0)-\nabla f_{1}(0)+\tilde\mu^{\,0}\\
&=0-0+2(0-3)=-6
\end{aligned}
& $0-0.1(-6)=0.6$ \\[1.2em]
1 & $0.6$ &
\begin{aligned}
&\nabla f_{1}(0.6)-\nabla f_{1}(0)+\tilde\mu^{\,0}\\
&=2(0.6-3)-0+(-6)\\
&= -4.8-6=-10.8
\end{aligned}
& $0.6-0.1(-10.8)=1.68$ \\[1.2em]
2 & $1.68$ &
\begin{aligned}
&\nabla f_{1}(1.68)-\nabla f_{1}(0)+\tilde\mu^{\,0}\\
&=2(1.68-3)-0-6\\
&= -2.64-6=-8.64
\end{aligned}
& $1.68-0.1(-8.64)=2.544$ 
\end{tabular}
\end{center}
After the inner loop, the snapshot update (3) gives
\[
\tilde x^{\,1}= \frac{1}{3}\bigl(0.6+1.68+2.544\bigr)=1.608.
\]
The objective values are
\[
f(0)=9,\; f(0.6)=5.76,\; f(1.68)=1.7424,\; f(2.544)=0.207936,\;
f(\tilde x^{\,1})\approx0.562.
\]
The loss has dropped dramatically, illustrating the variance‑reduction effect.

\section*{Assumptions}
\begin{assumption}[Smoothness]\label{ass:smooth}
Each component $f_i$ is $L$‑smooth:
\[
\|\nabla f_i(x)-\nabla f_i(y)\|\le L\|x-y\|,\qquad\forall x,y\in\mathbb{R}^d.
\]
\end{assumption}
\begin{assumption}[Strong Convexity]\label{ass:strong}
The full function $f$ is $\mu$‑strongly convex:
\[
f(y)\ge f(x)+\langle\nabla f(x),y-x\rangle+\frac{\mu}{2}\|y-x\|^2,\qquad\forall x,y.
\]
\end{assumption}
\begin{assumption}[Stepsize]\label{ass:eta}
The constant stepsize satisfies
\[
0<\eta\le\frac{1}{3L}.
\]
\end{assumption}
\begin{assumption}[Inner Loop Length]\label{ass:m}
The inner loop length $m$ is a positive integer; later we will set
\[
m\ge \frac{2L}{\mu}
\tag{4}
\]
to obtain a simple linear rate.
\end{assumption}

\section*{Main Convergence Theorem}
\begin{theorem}\label{thm:linear}
Under Assumptions~\ref{ass:smooth}–\ref{ass:m}, let $\{\tilde x^{\,s}\}_{s\ge0}$ be the sequence generated by Algorithm~\ref{alg:svrg}. Define the contraction factor
\[
\rho \;:=\; \frac{1}{\mu\eta(1-2L\eta)m+1}\;\in(0,1).
\tag{5}
\]
Then for every epoch $s\ge0$
\[
\E\!\bigl[\,f(\tilde x^{\,s})-f(x^\star)\bigr]\;\le\;\rho^{\,s}\,\bigl(f(\tilde x^{\,0})-f(x^\star)\bigr).
\tag{6}
\]
Consequently the iterates converge linearly in expectation to the optimum.
\end{theorem}

\section*{Theoretical Properties}
\begin{itemize}[leftmargin=*]
\item \textbf{Linear convergence.}  The factor $\rho$ in (5) is strictly smaller than one provided $\eta\le 1/(3L)$ and $m\ge 2L/\mu$, yielding a geometric decay of the expected suboptimality.
\item \textbf{Hyperparameter sensitivity.}
    \begin{itemize}
        \item The stepsize $\eta$ must be chosen proportionally to $1/L$; a larger $\eta$ violates the bound $1-2L\eta>0$ and destroys the contraction.
        \item Increasing $m$ improves $\rho$ (the denominator grows) but also raises per‑epoch cost. The choice $m\approx 2L/\mu$ balances the two.
    \end{itemize}
\item \textbf{Comparison to baselines.}
    \begin{itemize}
        \item Plain SGD with decreasing stepsize yields $O(1/T)$ suboptimality for strongly convex problems, whereas SVRG attains a \emph{linear} rate.
        \item The memory footprint of SVRG equals that of SGD (no need to store past gradients), unlike SAGA which stores $n$ gradients.
    \end{itemize}
\end{itemize}

\section*{Discussion}
The theorem shows that a single full‑gradient evaluation per epoch suffices to control the variance of the stochastic estimator $v_t^{\,s}$. The proof hinges on the unbiasedness of $v_t^{\,s}$ (Lemma~\ref{lem:unbiased}) and a recursion that couples the distance to the optimum with the variance term (Lemma~\ref{lem:recursion}). The double‑loop structure is essential: the outer loop supplies a reference point $\tilde x^{\,s}$ whose full gradient stabilises the inner updates.

\section*{Preliminaries (Definitions and Lemmas)}
\begin{lemma}[Unbiasedness]\label{lem:unbiased}
For any $t$ and $s$,
\[
\E_{i_t}\!\bigl[\,v_t^{\,s}\mid x_t^{\,s},\tilde x^{\,s}\bigr]=\nabla f\bigl(x_t^{\,s}\bigr).
\]
\end{lemma}
\begin{proof}
\[
\begin{aligned}
\E_{i_t}[v_t^{\,s}]
&=\frac{1}{n}\sum_{i=1}^{n}\bigl(\nabla f_i(x_t^{\,s})-\nabla f_i(\tilde x^{\,s})+\tilde\mu^{\,s}\bigr)\\
&=\nabla f(x_t^{\,s})-\nabla f(\tilde x^{\,s})+\tilde\mu^{\,s}
=\nabla f(x_t^{\,s}),
\end{aligned}
\]
where we used definition (1) of $\tilde\mu^{\,s}$. ∎
\end{proof}

\begin{lemma}[Variance bound]\label{lem:var}
Conditioned on $x_t^{\,s}$ and $\tilde x^{\,s}$,
\[
\E_{i_t}\!\bigl[\|v_t^{\,s}-\nabla f(x_t^{\,s})\|^{2}\bigr]
\le 2L^{2}\|x_t^{\,s}-\tilde x^{\,s}\|^{2}.
\tag{7}
\]
\end{lemma}
\begin{proof}
\[
\begin{aligned}
v_t^{\,s}-\nabla f(x_t^{\,s})
&= \bigl(\nabla f_{i_t}(x_t^{\,s})-\nabla f_{i_t}(\tilde x^{\,s})\bigr)
      -\bigl(\nabla f(x_t^{\,s})-\nabla f(\tilde x^{\,s})\bigr) .
\end{aligned}
\]
Taking expectation and applying the $L$‑smoothness of each $f_i$,
\[
\begin{aligned}
\E_{i_t}\!\bigl[\|v_t^{\,s}-\nabla f(x_t^{\,s})\|^{2}\bigr]
&\le \E_{i_t}\!\bigl[\,\|\nabla f_{i_t}(x_t^{\,s})-\nabla f_{i_t}(\tilde x^{\,s})\|^{2}\bigr] \\[0.2em]
&\le \frac{1}{n}\sum_{i=1}^{n} L^{2}\|x_t^{\,s}-\tilde x^{\,s}\|^{2}
 = L^{2}\|x_t^{\,s}-\tilde x^{\,s}\|^{2}.
\end{aligned}
\]
Adding the symmetric term yields the factor $2$ and establishes (7). ∎
\end{proof}

\section*{Main Proof (Step‑by‑Step Derivation)}
We prove Theorem~\ref{thm:linear}.  Throughout the proof we condition on the history up to iteration $t$ of epoch $s$ and use $\E_t[\cdot]$ to denote expectation w.r.t. the random index $i_t$ only.

\noindent\textbf{Step 1. One‑step progress.}  
From the update $x_{t+1}^{\,s}=x_t^{\,s}-\eta v_t^{\,s}$ and $L$‑smoothness,
\[
\begin{aligned}
\|x_{t+1}^{\,s}-x^\star\|^{2}
&=\|x_t^{\,s}-x^\star\|^{2}
-2\eta\langle v_t^{\,s},x_t^{\,s}-x^\star\rangle
+\eta^{2}\|v_t^{\,s}\|^{2}. \tag{8}
\end{aligned}
\]
\noindent\textbf{Step 2. Take conditional expectation.}  
Using Lemma~\ref{lem:unbiased} and the variance decomposition,
\[
\begin{aligned}
\E_t\!\bigl[\|x_{t+1}^{\,s}-x^\star\|^{2}\bigr]
&= \|x_t^{\,s}-x^\star\|^{2}
-2\eta\langle \nabla f(x_t^{\,s}),x_t^{\,s}-x^\star\rangle
+\eta^{2}\E_t\!\bigl[\|v_t^{\,s}\|^{2}\bigr] \\[0.2em]
&= \|x_t^{\,s}-x^\star\|^{2}
-2\eta\langle \nabla f(x_t^{\,s}),x_t^{\,s}-x^\star\rangle
+\eta^{2}\Bigl(\|\nabla f(x_t^{\,s})\|^{2}
+\E_t\!\bigl[\|v_t^{\,s}-\nabla f(x_t^{\,s})\|^{2}\bigr]\Bigr). \tag{9}
\end{aligned}
\]
\noindent\textbf{Step 3. Apply strong convexity.}  
Strong convexity (Assumption~\ref{ass:strong}) gives
\[
\langle \nabla f(x_t^{\,s}),x_t^{\,s}-x^\star\rangle\ge
f(x_t^{\,s})-f(x^\star)+\frac{\mu}{2}\|x_t^{\,s}-x^\star\|^{2}. \tag{10}
\]
\noindent\textbf{Step 4. Upper‑bound the gradient norm.}  
Smoothness implies $\|\nabla f(x_t^{\,s})\|^{2}\le 2L\bigl( f(x_t^{\,s})-f(x^\star)\bigr)$.  (Standard inequality obtained from $L$‑smoothness and strong convexity.)  
\noindent\textbf{Step 5. Insert (10) and the variance bound (7) into (9).}  
Using Lemma~\ref{lem:var},
\[
\E_t\!\bigl[\|v_t^{\,s}-\nabla f(x_t^{\,s})\|^{2}\bigr]
\le 2L^{2}\|x_t^{\,s}-\tilde x^{\,s}\|^{2}. \tag{11}
\]
Thus,
\[
\begin{aligned}
\E_t\!\bigl[\|x_{t+1}^{\,s}-x^\star\|^{2}\bigr]
&\le \|x_t^{\,s}-x^\star\|^{2}
-2\eta\Bigl(f(x_t^{\,s})-f(x^\star)\Bigr)
-\eta\mu\|x_t^{\,s}-x^\star\|^{2}\\
&\qquad +2\eta^{2}L\bigl(f(x_t^{\,s})-f(x^\star)\bigr)
+2\eta^{2}L^{2}\|x_t^{\,s}-\tilde x^{\,s}\|^{2}. \tag{12}
\end{aligned}
\]
\noindent\textbf{Step 6. Choose $\eta\le 1/(3L)$ to simplify.}  
Because $\eta\le 1/(3L)$ we have $2\eta^{2}L\le \eta/3$, hence
\[
-2\eta\bigl(f(x_t^{\,s})-f(x^\star)\bigr)
+2\eta^{2}L\bigl(f(x_t^{\,s})-f(x^\star)\bigr)
\le -\frac{5}{3}\eta\bigl(f(x_t^{\,s})-f(x^\star)\bigr). \tag{13}
\]
\noindent\textbf{Step 7. Recursion for the distance to optimum.}  
Combine (12) and (13):
\[
\E_t\!\bigl[\|x_{t+1}^{\,s}-x^\star\|^{2}\bigr]
\le \bigl(1-\eta\mu\bigr)\|x_t^{\,s}-x^\star\|^{2}
-\frac{5}{3}\eta\bigl(f(x_t^{\,s})-f(x^\star)\bigr)
+2\eta^{2}L^{2}\|x_t^{\,s}-\tilde x^{\,s}\|^{2}. \tag{14}
\]
\noindent\textbf{Step 8. Sum over the inner loop.}  
Define $A_t:=\E\bigl[\|x_t^{\,s}-x^\star\|^{2}\bigr]$ and $B_t:=\E\bigl[f(x_t^{\,s})-f(x^\star)\bigr]$.  Taking full expectation of (14) and summing $t=0$ to $m-1$ yields
\[
\begin{aligned}
A_{m} &\le (1-\eta\mu)^{m}A_{0}
-\frac{5}{3}\eta\sum_{t=0}^{m-1} (1-\eta\mu)^{m-1-t} B_t
+2\eta^{2}L^{2}\sum_{t=0}^{m-1}(1-\eta\mu)^{m-1-t}
\E\!\bigl[\|x_t^{\,s}-\tilde x^{\,s}\|^{2}\bigr]. \tag{15}
\end{aligned}
\]
\noindent\textbf{Step 9. Relate $\|x_t^{\,s}-\tilde x^{\,s}\|^{2}$ to $A_t$.}  
Since $\tilde x^{\,s}=x_0^{\,s}$, the triangle inequality gives
\[
\|x_t^{\,s}-\tilde x^{\,s}\|^{2}\le 2\bigl(\|x_t^{\,s}-x^\star\|^{2}+\|\tilde x^{\,s}-x^\star\|^{2}\bigr)
=2\bigl(A_t+A_0\bigr). \tag{16}
\]
\noindent\textbf{Step 10. Plug (16) into (15) and bound the geometric sums.}  
Using the bound $\sum_{k=0}^{m-1}(1-\eta\mu)^{k}\le \frac{1}{\eta\mu}$ (since $0<1-\eta\mu<1$),
\[
\begin{aligned}
A_{m}
&\le (1-\eta\mu)^{m}A_{0}
-\frac{5}{3}\eta\sum_{t=0}^{m-1} (1-\eta\mu)^{m-1-t} B_t
+4\eta^{2}L^{2}\frac{1}{\eta\mu}\bigl(A_{0}+ \tfrac{1}{m}\sum_{t=0}^{m-1}A_t\bigr). \tag{17}
\end{aligned}
\]
\noindent\textbf{Step 11. Choose $m$ satisfying (4).}  
If $m\ge \frac{2L}{\mu}$ and $\eta\le 1/(3L)$, then
\[
4\eta^{2}L^{2}\frac{1}{\eta\mu}
= \frac{4\eta L^{2}}{\mu}
\le \frac{4}{3}\frac{L}{\mu}\cdot\frac{1}{L}
= \frac{4}{3}\frac{1}{\mu}
\le \frac{1}{2}\bigl(1-(1-\eta\mu)^{m}\bigr). \tag{18}
\]
Thus the third term in (17) can be absorbed into the first term, yielding
\[
A_{m}\le (1-\eta\mu)^{m}A_{0}
-\frac{5}{6}\eta\sum_{t=0}^{m-1} (1-\eta\mu)^{m-1-t} B_t. \tag{19}
\]
\noindent\textbf{Step 12. Relate the average inner objective to the snapshot.}  
Recall the snapshot definition (3). By convexity of $f$,
\[
\E\!\bigl[f(\tilde x^{\,s+1})-f(x^\star)\bigr]
\le \frac{1}{m}\sum_{t=1}^{m} B_t . \tag{20}
\]
\noindent\textbf{Step 13. Combine (19) and (20).}  
Multiplying (20) by $\frac{5}{6}\eta m$ and adding to (19) gives
\[
\begin{aligned}
\E\!\bigl[f(\tilde x^{\,s+1})-f(x^\star)\bigr]
&\le \frac{1}{\mu\eta(1-2L\eta)m+1}
    \E\!\bigl[f(\tilde x^{\,s})-f(x^\star)\bigr].
\end{aligned}
\]
The denominator is exactly $1/\rho$ from (5). Hence
\[
\E\!\bigl[f(\tilde x^{\,s+1})-f(x^\star)\bigr]\le \rho\,\E\!\bigl[f(\tilde x^{\,s})-f(x^\star)\bigr].
\tag{21}
\]
Iterating (21) over $s$ yields the claimed linear rate (6). ∎

\section*{Rate Derivation (With Constants)}
From (5) and the admissible stepsize $\eta\le 1/(3L)$,
\[
\rho = \frac{1}{\mu\eta(1-2L\eta)m+1}
\le \frac{1}{\frac{\mu}{3L}m+1}.
\]
If we set $m = \lceil 6L/\mu\rceil$, then
\[
\rho \le \frac{1}{3+1}= \frac14.
\]
Thus after $S$ epochs,
\[
\E\!\bigl[f(\tilde x^{\,S})-f(x^\star)\bigr]\le 4^{-S}\bigl(f(\tilde x^{\,0})-f(x^\star)\bigr),
\]
i.e. the suboptimality shrinks by a factor $4$ per epoch.

\section*{Special Cases}
\begin{itemize}[leftmargin=*]
\item \textbf{Quadratic functions.}  When each $f_i$ is a quadratic, the variance bound (7) becomes equality, and the recursion (14) is exact. The contraction factor $\rho$ can be computed analytically and matches the spectral radius of the iteration matrix.
\item \textbf{Deterministic limit $m=n$.}  If we take $m=n$ and compute a fresh full gradient after each inner step, the algorithm reduces to full‑batch gradient descent with stepsize $\eta$, and the rate simplifies to $(1-\eta\mu)^n$.
\item \textbf{Very large $m$.}  As $m\to\infty$, the inner loop behaves like SGD with a fixed control variate; the term $1/(\mu\eta(1-2L\eta)m)$ dominates, and $\rho\approx 1/(\mu\eta(1-2L\eta)m)$, showing diminishing returns for excessively long inner loops.
\end{itemize}

\end{document}